\section{From the height error to a one-sided prime-sum estimate}
\label{sec:primecomparison}

Fix a scale \(x\geq3\) and a tolerance \(\varepsilon>0\).
The comparison in this section takes place separately at each scale.
The finite groups, weights, sets and auxiliary parameters may vary
with \(x\). We display the constants that have to be fixed before the
scale is chosen.

\subsection{The finite Fourier identity, including ramified primes}

Let \(H\) be a finite abelian group of order \(n\), written additively,
and put \(E=(\mathbb Z/2)^2\times H\),
\(\Phi=\{(0,v,h)\}\subset E\), \(f=|\Phi|=2n\).
Elements denoted \(t\in H\) are identified with \((0,0,t)\) in \(E\).
Let \(\kappa=(1,1,0)\), \(e=(1,0,0)\),
\(\xi(v_1,v_2,h)=(-1)^{v_1}\), and
\(\rho_t(g)=\kappa+g+t\). The map \(\rho_t\) carries the
complement of \(\Phi\) bijectively onto \(\Phi\).

The arithmetic input at a prime \(p\) is an inertia subgroup \(I_p\)
of \(E\) and a Frobenius representative \(F_p\in E\). Only its
coset modulo \(I_p\) enters the sums. For a finite order character
\(\psi\in\widehat E\), its value at a prime power in this construction is
\[
 v_\psi(p^i)=\frac1{|I_p|}\sum_{h\in I_p}\psi(iF_p+h)
 \qquad(i\geq1).
\]
Indeed restriction of \(\psi\) to inertia is either trivial, in
which case the sum is \(\psi(iF_p)\), or has zero mean, in which
case the local character contributes zero. The host construction
identifies these values with those of its associated primitive
Dirichlet character. Orthogonality now gives an exact formula,
\begin{equation}
 \frac1f\sum_{\psi\in\widehat E}\overline{\psi(g)}v_\psi(p^i)
 =\frac{2}{|I_p|}{\bf1}_{\,g-iF_p\in I_p}.
 \label{eq:pc-inertia}
\end{equation}
To see this, insert the average defining \(v_\psi\). The inner
character sum is \(|E|=2f\) when \(iF_p+h=g\) and is zero otherwise.
There is at most one such \(h\).

Here the host provides a finite set \(S\) containing \(2,q_1\) and
a prime \(d\notin S\). At \(p\in S\) it gives \(I_p\subset\Phi\)
and \(F_p\in\Phi\), at \(d\) it gives \(I_d=\{0,e\}\), and at all
other primes it gives \(I_p=0\). Let \(w_\rho>0\) for \(\rho\in\Phi\),
\(A>0\), and \(N\geq A/w_\rho\) for every \(\rho\). Write
\(r_{m,g}=(A/m-w_{\rho_t(g)})_+\), where \(z_+=\max(z,0)\).
The arithmetic prime contribution for this \(A,t,w\) equals
\begin{align*}
 P(A,t,w)={}&\sum_{\substack{p\leq N,\ p\notin S\cup\{d\}\\ F_p\notin\Phi}}
  2\log p\sum_{\substack{1\leq i\leq N\\iF_p\notin\Phi}}r_{p^i,iF_p}\\
 &+\log d\sum_{1\leq i\leq N}
       \sum_{\substack{g\notin\Phi\\g-iF_d\in I_d}}r_{d^i,g}.
\end{align*}
The sums indexed by \(p\) in this section range over primes. When
\(p^i>N\) the ramp vanishes, so the second line is zero if \(d>N\).
In view of \eqref{eq:pc-inertia}, the same real number is
\begin{equation}
 P(A,t,w)=\frac1f\sum_{g\notin\Phi}\sum_{\psi\in\widehat E}
  \overline{\psi(g)}
  \sum_{m=1}^N\Lambda(m)v_\psi(m)(A/m-w_{\rho_t(g)})_+ .
 \label{eq:pc-prime}
\end{equation}
At \(p\in S\) the entire coset is in \(\Phi\). At an unramified
prime the surviving term has coefficient two and \(g=iF_p\).
At \(d\) the surviving coset terms have coefficient one. These are
all the terms, since \(\Lambda(m)\) vanishes away from prime powers.

Let \(G_\psi(y)\) be the ramp with coefficients \(\Lambda(m)v_\psi(m)\),
centered by \(y\ell_\psi\) for \(\psi\ne1\), where \(\ell_\psi\)
is the logarithmic derivative of the associated primitive \(L\)-function
at one. For \(\psi=1\) the centering is
\(-y\log y+(1+\gamma)y\). These formulas also define the functions
for \(0<y<1\) by taking the ramp sum to be empty. Put
\[
 M(A,w)=A\left(\log A-1-\gamma+\ell_\xi
                         -\frac1f\sum_{\rho\in\Phi}\log w_\rho\right).
\]
Then
\begin{equation}
 P(A,t,w)-M(A,w)
   =\frac1f\sum_{g\notin\Phi}\sum_{\psi\in\widehat E}
     \overline{\psi(g)}\,w_{\rho_t(g)}G_\psi(A/w_{\rho_t(g)}).
 \label{eq:pc-balance}
\end{equation}
For \(\psi\ne1\), multiplying the centered ramp by \(w\)
adds \(A\ell_\psi\) to the inner prime sum of \eqref{eq:pc-prime}.
For \(\psi=1\) it adds \(-A\log A+A\log w+(1+\gamma)A\).
The sum of \(\overline{\psi(g)}\) over \(g\notin\Phi\) is \(f\)
at \(1\), \(-f\) at \(\xi\), and zero otherwise, by character
orthogonality on the index-two subgroup. The bijection \(\rho_t\)
turns the logarithmic sum into the sum over \(\Phi\) in \(M\).
Moving the two corrections to the left gives \eqref{eq:pc-balance}.

\subsection{The weights and the surviving characters}

Let \(H_0\subset H\) be a subgroup, \(T\subset H\), \(T_2\subset H_0\),
\(n_T=|T|\), \(n_2=|T_2|\) and \(\Delta=n_2/n\).
For \(u\in H_0\), the thick positions are precisely
\(\rho=(0,1,s+u)\), \(s\in T_2\).
For each \(1\leq j\leq m\) take positive weights \(w_{u,j,\rho}\).
Take \(b,b',X'>0,\ \eta,\vartheta\geq0\) such that
\[
 b'X'\leq xb,\qquad b\leq w_{u,j,\rho}\leq b+\eta+1
 \ \ (\rho\text{ thick}),\qquad
 w_{u,j,\rho}\leq2b'\ \ (\rho\text{ otherwise}).
\]
At \(t\in T\) put \(A=a_t\), \(y_t=a_t/b\), and require
\(|y_t-x|\leq1\) and \(a_t/w_{u,j,\rho}\leq N\).
The arithmetic estimates on the thin scales and the window costs are
\begin{equation}
 |G_\psi(y)|\leq y\vartheta\quad(y\geq X'/3),\qquad
 2f(x+1)\vartheta\leq x^{-1},\qquad
 6f(x+1)(\eta+1)\leq b/x .
 \label{eq:pc-cost}
\end{equation}

The thick part after replacing each thick weight by \(b\) is
\[
 T_u(t)=\frac1f\sum_{s\in T_2}\sum_{\psi\in\widehat E}
          \psi\bigl(-(e+s+u-t)\bigr)bG_\psi(y_t).
\]
Indeed \(\rho_t(g)=(0,1,s+u)\) if and only if \(g=e+s+u-t\).
We claim
\begin{equation}
 \Re T_u(t)\leq P(a_t,t,w_{u,j})-\Re M(a_t,w_{u,j})+2b/x .
 \label{eq:pc-individual}
\end{equation}
For the thick terms the inequality used is
\[
 |wG_\psi(A/w)-bG_\psi(A/b)|
                  \leq3A\,|w-b|/\min(w,b).
\]
For a nontrivial \(\psi\) the two linear centering terms cancel.
The ramps are \(1\)-Lipschitz in \(w\) and vanish for
\(m>A/\min(w,b)\). The classical Chebyshev bound
\(\sum_{m\leq z}\Lambda(m)\leq3z\) proves the claim.
For the trivial \(\psi\), suppose first \(b\leq w\). The decrease
of the uncentered ramp sum lies between \(0\) and \(3A(w-b)/b\);
the change from the centering has the opposite sign and size
\(A\log(w/b)\leq A(w-b)/b\). The same bound therefore holds.
Interchanging \(w,b\) covers the other order. Thus the replacement
cost at a thick term is at most \(3(x+1)(\eta+1)\).
At a thin term \(w>0\), \(w\leq2b'\) and \(A\geq(x-1)b\).
Since \(x\geq3\), these give \(A/w\geq X'/3\).
Its absolute contribution is at most \(A\vartheta\leq(x+1)b\vartheta\).
In \eqref{eq:pc-balance} there are \(f\) choices for \(g\), \(2f\)
characters and the factor \(1/f\). Hence its change in absolute value
is bounded by \(2f[3(x+1)(\eta+1)+(x+1)b\vartheta]\leq2b/x\)
by \eqref{eq:pc-cost}. This proves \eqref{eq:pc-individual}.

Let \(\mathcal A\subset\widehat E\) consist of the characters trivial
on \(H_0\). For each \(s\in T_2\), orthogonality on \(H_0\) gives
\[
 \sum_{u\in H_0}\psi(-(e+s+u-t))=
 \begin{cases}|H_0|\psi(e+t),&\psi\in\mathcal A,\\0,&\psi\notin\mathcal A.\end{cases}
\]
Here \(e=-e\) and \(\psi(s)=1\) in the first case. It follows that
\[
 \frac1{|H_0|}\sum_{u\in H_0}T_u(t)
      =b\,\frac{\Delta}{2}\sum_{\psi\in\mathcal A}\psi(e+t)G_\psi(y_t).
\]
Write \(E_{\rm top}\) for a common upper bound, for every \(u\in H_0\),
on \(\sum_{t\in T}\sum_{j=1}^m [P(a_t,t,w_{u,j})-\Re M(a_t,w_{u,j})]\).
Average \eqref{eq:pc-individual} over \(u\) and sum over \(t,j\).
The result is
\begin{equation}
 mb\frac{\Delta}{2}\sum_{t\in T}\sum_{\psi\in\mathcal A}
       \Re[\psi(e+t)G_\psi(y_t)]
              \leq E_{\rm top}+2n_Tmb/x .
 \label{eq:pc-top}
\end{equation}

\subsection{Comparison with a fixed family of primitive characters}

Let \(K\subset\mathcal A\) contain \(1\). Suppose that \(\psi(e)=1\)
on \(K\). Set \(c_\psi=1\) if \(\psi\) is trivial on \(H\), and
\(c_\psi=c\) otherwise. Fix \(|c|\leq1\) and constants \(M_0,K_0,C_0\)
valid throughout the windows, with
\begin{gather*}
 |\mathcal A|\leq M_0,\qquad
 \left|\sum_{t\in T}\psi(t)-c_\psi n_T\right|
                  \leq M_0x^{-1+\varepsilon}n_T\quad(\psi\in K),\\
 |\ell_\psi|\leq K_0\quad(\psi\in K,\ \psi\ne1),\qquad
 |\ell_\psi|\leq C_0\log x\quad(\psi\in\mathcal A\setminus K).
\end{gather*}
For \(\psi\) trivial on \(H\) the first sum equals \(n_T\) exactly.
Define
\[
 G_K=\sum_{\psi\in K}c_\psi\Re G_\psi(x),\qquad
 \Sigma=\sum_{\psi\in\mathcal A\setminus K}|G_\psi(x)|,\qquad
 V=G_K-\Sigma .
\]
For \(x\) sufficiently large depending only on the fixed constants
and \(\varepsilon\), we have
\begin{equation}
 n_T(V-4x^{1/2+\varepsilon})\leq
  \sum_{t\in T}\sum_{\psi\in\mathcal A}\Re[\psi(e+t)G_\psi(y_t)] .
 \label{eq:pc-lower}
\end{equation}

For the proof, the ramp is locally Lipschitz. Off the integers the
derivative of the principal centered ramp is
\(\sum_{m\leq y}\Lambda(m)/m-\log y+\gamma\).
The elementary Mertens estimate
\(\sum_{m\leq y}\Lambda(m)/m=\log y+O(1)\) therefore gives an absolute
bound \(D\) for this derivative on \(y\geq1\). Continuity at integers
extends the Lipschitz estimate across them. For a nontrivial \(\psi\),
the absolute derivative is at most \(\log y+D+|\ell_\psi|\) by the
same estimate, increasing \(D\geq0\) once if necessary. Since
\(G_\psi(1)=\ell_\psi\) for \(\psi\ne1\), we obtain
\begin{align*}
 |G_\psi(y_t)-G_\psi(x)|&\leq\log(x+1)+D+|\ell_\psi|,\\
 |G_\psi(x)|&\leq x(\log x+D+|\ell_\psi|).
\end{align*}
For \(1\ne\psi\in K\), subtract \(n_Tc_\psi G_\psi(x)\) from its
sum in \eqref{eq:pc-lower} before taking real parts. The difference is
\[
  \sum_{t\in T}\psi(t)[G_\psi(y_t)-G_\psi(x)]
          +\left(\sum_{t\in T}\psi(t)-c_\psi n_T\right)G_\psi(x).
\]
In absolute value this is at most \(n_T\) times
\(\log(x+1)+D+|K_0|+|M_0|x^\varepsilon(\log x+D+|K_0|)\).
For \(\psi=1\) the lower error per term is at most \(D\).
For \(\psi\in\mathcal A\setminus K\) the real part of each term is
bounded below by \(-|G_\psi(x)|-\log(x+1)-D-|C_0|\log x\).
Thus all three errors per term are at most \(A_1x^{1/4+\varepsilon}\), where
\[
 A_1=5+D+|K_0|+|M_0|(4+D+|K_0|)+4|C_0|.
\]
Indeed \(\log x\leq4x^{1/4}\), \(\log(x+1)\leq5x^{1/4}\) for \(x\geq3\).
Enlarge the threshold so that \(|M_0|A_1\leq x^{1/4}\). The total
per \(t\) is then at most \(x^{1/2+\varepsilon}\), which gives the
stated bound with the source's constant four.

\subsection{Accounting for the size error}

The source supplies a geometric upper bound for \(E_{\rm top}\) at
this same \(x\). In numerical form, with \(C\geq0\), \(a_E\geq0\), it is
\begin{align}
 E_{\rm top}\leq x^{C\varepsilon}\bigg[
 n_T a_E(1+m^2/k)
 +m\left(2nmbx^\varepsilon+
       \frac{n_2 k b x^\varepsilon}{n(n x^{-\varepsilon})}\right)\bigg].
 \label{eq:pc-geometry}
\end{align}
The size hypotheses that enter division by the main coefficient are
\[
 \begin{gathered}
 0\leq a_E\leq2xb,\quad n_T\geq nx^{-\varepsilon/8}/2,
 \quad x^{-\varepsilon}\leq\Delta\leq1,\\
 x^{1/2-\varepsilon}\leq m\leq x^{1/2+\varepsilon/2},\qquad
 xm\leq k\leq n^2x^{1/2-\varepsilon},\qquad
 26\leq x^{3\varepsilon/8}.
 \end{gathered}
\]
Let \(Q=n_Tmb>0,\ Y=x^{1/2+13\varepsilon/8},\ P_0=x^{C\varepsilon}\).
Divide the four summands inside the bracket by \(Q\).
The first is \(a_E/(mb)\leq2x/m\leq2x^{1/2+\varepsilon}\leq2Y\).
The second is \(a_E m/(bk)\leq2xm/k\leq2\leq2Y\).
The third is \(2nm x^\varepsilon/n_T
      \leq4x^{1/2+13\varepsilon/8}=4Y\).
The fourth is \(n_2 k x^{2\varepsilon}/(n_Tn^2)
      \leq2x^{1/2+9\varepsilon/8}\leq2Y\).
Hence \(E_{\rm top}\leq10P_0YQ\). Combining
\eqref{eq:pc-lower} and \eqref{eq:pc-top}, dividing by \(Q\)
and using \(2/x\leq1\leq P_0Y\), gives
\[
 \Delta(V-4x^{1/2+\varepsilon})\leq22P_0Y .
\]
As \(\Delta^{-1}\leq x^\varepsilon\) and \(C\geq0\), it follows that
\begin{equation}
 V\leq26x^{1/2+(C+21/8)\varepsilon}
       \leq x^{1/2+(C+3)\varepsilon}.
 \label{eq:pc-core}
\end{equation}

For the cyclic level of order \(o\) used to extract a prescribed
primitive character, the host parameters are
\[
 c=(3o+1)^{-1},\qquad M_0=4o,\qquad
 K_0=\sum_{\chi\ \mathrm{retained}} |L'(1,\chi)/L(1,\chi)|
\]
with the formal value for the principal component in \(K_0\) understood
as the library's totalized logarithmic derivative. Equivalently one
may increase any bound for the nonprincipal summands by a fixed
nonnegative constant, since it is used only on \(\psi\ne1\).
The character group \(\mathcal A\) is the dual of \(E/H_0\),
so \(|\mathcal A|=4[H:H_0]=4o\).
The arithmetic host identifies the retained characters bijectively
with the primitive family of the level and identifies their ramps
and weights, hence \(G_K=G_{\mathcal L}(x)\).
The host's bounded-character estimate gives
\(|G_\psi(x)|\leq x^{1/2+\varepsilon}\) on
\(\mathcal A\setminus K\), whence \(\Sigma\leq4o x^{1/2+\varepsilon}\).
The height error exponent at this scale is at most \(e_D+e_L+3\)
with \(e_D+e_L\geq0\). Thus \eqref{eq:pc-geometry} may be used with
\(C=e_D+e_L+3\). Adding \(\Sigma\) in \eqref{eq:pc-core} and requiring
\(1+4o\leq x^\varepsilon\) gives
\[
 G_{\mathcal L}(x)\leq x^{1/2+(e_D+e_L+7)\varepsilon}.
\]
Every large-scale threshold here depends on the fixed level, its
error constants and \(\varepsilon\). Consequently, if the arithmetic
inputs hold at every point of one interval above those thresholds,
the conclusion holds at every point of that same interval.

\paragraph{Correspondence and inputs.}
The Fourier calculation translates \texttt{ch\_\allowbreak{}balance}.
The projection, replacement of the weights, and character comparison
translate \texttt{Character\allowbreak{}Balance\allowbreak{}Moves.\allowbreak{}sum\_\allowbreak{}Th}, \texttt{th\_\allowbreak{}le},
\texttt{one\_\allowbreak{}char}, \texttt{core} and \texttt{ch\_\allowbreak{}core}.
The four-term calculation translates \texttt{proposition\_\allowbreak{}size\_\allowbreak{}bound}.
The conversion at a fixed level translates \texttt{ch\_\allowbreak{}top},
\texttt{gm\allowbreak{}A\_\allowbreak{}eq}, \texttt{sig\allowbreak{}A\_\allowbreak{}le} and \texttt{core\_\allowbreak{}gm}.
The upstream inputs used here are the host with the specified
Frobenius and inertia data, the common-window character and weight
estimates, the inequalities on their sizes, and the geometric
inequality \eqref{eq:pc-geometry}. In \texttt{ch\_\allowbreak{}top} this last
inequality follows from the source's unitwise height bound
\texttt{each\_\allowbreak{}unit\_\allowbreak{}bound\_\allowbreak{}statement} and error comparison
\texttt{err\allowbreak{}X5at\_\allowbreak{}le} after substituting the upper bounds for its
two coefficients \(B_0,N_c\).
